% chapter7.tex
\section{总结与展望}

\subsection{研究工作总结}
本课题以可变形物体实时触摸建模和触觉交互仿真为研究核心，以Sapienza/INSA大学的DeformableSimulation开源项目（原用于腹腔手术触觉训练）为基础进行改造优化，系统研究了柔性体物理仿真、触觉渲染、逆运动学求解、多线程同步、视触融合等关键技术，最后设计并实现了基于MuJoCo 3.x物理引擎的可变形物体实时触摸交互仿真系统，完成了系统架构设计、硬件依赖解耦、软件开发、算法实现、场景优化、实验验证的全过程，形成了一个完整的感知、计算、反馈、评价的闭环交互体系，达到了预期的研究目的。

按照分层模块化的架构设计思想，创建起交互层、GUI、主循环、引擎层、触觉反馈算法层这五层结构，符合Python 3.10+全栈开发环境的要求，把mujoco.viewer可视化、pynput键盘交互、numpy数值计算这些核心技能整合起来，解决之前XML兼容性的问题，加入中文界面说明，改善场景资产的可利用程度，达成双手加腹腔器官（肝脏，肾脏等）的演示场景，允许用户用键盘操控虚拟手进行软组织接触、按压等互动操作，显示即时形变反应。

核心技术上使用MuJoCo的flexcomp柔性组件对腹腔器官进行网格化建模，把器官离散成质点网格，用edge equality约束和阻尼连接建立mass-spring FEM网络，用solimp/solref参数定义接触特性，实现不同器官独立物理参数（质量、刚度、阻尼、摩擦）的精确设置；改进虚拟手运动控制，用改良阻尼最小二乘法求解逆运动学，结合手部mocap驱动和手指IK计算，保证虚拟手动作流畅自然；抽象并实现独立的触觉反馈算法模块（haptic\_feedback.py），包含力信号滤波、自适应力计算、纹理映射、多物体接触检测四个主要功能，完全不依赖硬件，从算法层面实现了完整的触觉建模。

项目改造及适配方面主要完成三项工作：一是解耦原项目对于VR设备和WEART触觉手套的强依赖，设计并实现键盘交互模式，适应无硬件实验环境；二是把原项目里耦合到硬件调用的力反馈逻辑，抽象成独立的触觉反馈算法模块，达成纹理解耦和算法通用化，给以后接入任意触觉硬件赋予支撑；三是改善场景可用性，解决XML文件兼容性问题，改进组合场景设置，保证系统正常运作。实验结果表明，系统可以稳定地完成可变形器官的高精度建模和实时交互，触觉反馈算法可以很好地消除仿真抖动，还原真实的触感特征，完全满足了“模拟器内采集接触力和视觉形变”的设计要求。

\subsection{研究主要创新点}
结合项目改造优化工作和技术实现细节，本文主要的创新点有系统适配、算法设计、模块架构三个方面，均为本人独立完成的工作，具体如下所示。

\begin{enumerate}
    \item \textbf{适配无硬件实验环境，实现交互模式降级优化}：针对原项目需要VR头显和触觉手套的限制，解耦硬件依赖，设计pynput跨平台键盘监听交互模式，完成虚拟手“瞬移式”驱动适配，优化FEM求解器参数，尽量减少降级交互造成的非物理穿透和求解器发散问题，提高系统在普通实验环境下的可用性和稳定性。
    
    \item \textbf{抽象触觉反馈算法模块，实现算法和硬件解耦}：把原项目里耦合在WEART硬件调用上的力反馈逻辑，提取出来并重新构造为一个独立的触觉反馈算法层，该层包含ForceFilter滑动平均滤波、AdaptiveForceCalculator自适应力计算、TextureMapper纹理映射这三个主要模块，可以对不同的组织触觉特征进行独立的映射以及动态调节，摆脱硬件绑定的束缚，以后可以自由地接入各种触觉反馈设备，提高算法的通用性和可扩展性。
    
    \item \textbf{改善场景及系统可用性，提高项目落地性}：修复原项目XML文件兼容性问题，完善双手和腹腔器官组合场景的配置，增加中文化GUI操作说明，改进主循环（simulator.py）的输入、物理步进、渲染、反馈流程，提高系统的运行稳定性，补充完整的场景资产管理，明确hands.xml、phantom.xml等核心文件的功能和用途，为以后系统的维护和功能扩展提供方便。
    
    \item \textbf{完善系统评价体系，加强技术可验证性}：根据系统实际运行情况补充键盘交互模式、触觉算法有效性验证内容，从物理仿真精度、交互流畅度、算法稳定性三个方面验证系统在无硬件环境下的可行性，形成改造、实现、验证的完整技术闭环，给同类项目硬件解耦和算法优化提供参考。
\end{enumerate}

\subsection{研究不足与限制}
根据系统实际运行情况和答辩预判结果可知，本研究在硬件依赖、交互精度、求解器稳定性这三个方面存在一定的不足和限制，和前面第六章的局限性分析是一致的，如下所示。

\begin{enumerate}
    \item \textbf{硬件受限，缺少完整的回路人机交互验证}：由于实验条件所限，没有接入WEART TouchDIVER触觉手套和Oculus Rift S VR头显，不能得到真实的力反馈输出以及连续的手部追踪，只能进行算法层面的触觉建模和视觉形变仿真，系统的真正交互体验及临床应用价值还没有得到充分的验证。
    
    \item \textbf{交互精度不够，降级交互存在固有的缺陷}：为了适应没有硬件环境的键盘离散输入模式，使虚拟手呈现出“瞬移式”的驱动方式，与原项目VR头显连续mocap追踪的设计初衷存在差异，容易在与软体接触时出现非物理穿透，造成FEM求解器发散，虚拟手飞出视野、器官模型“爆炸”等异常情况，影响交互的稳定性和真实性。
    
    \item \textbf{算法和仿真精度还有待提高}：触觉反馈算法中纹理映射仍然使用预先定义好的纹理库，没有根据真实的生物组织摩擦特性动态生成，精细度与真实的材质有差距；MuJoCo的FEM求解器采用简化质点-弹簧模型，在极端大变形的情况下数值精度不高，不能完全反映腹腔器官复杂的力学响应。
    
    \item \textbf{系统功能单一}：只能进行单手三指的键盘交互，没有实现双手协同、多指精细控制等复杂的操作，场景规模小，可以交互的物体种类只有腹腔核心器官，在多物体紧密聚集的情况下接触检测效率较低。
\end{enumerate}

\subsection{未来研究与应用展望}
根据本研究的不足之处，结合原项目临床应用定位以及目前的技术发展走向，未来将会从硬件接入、算法优化、功能拓展这三个主要方面展开深入探究，使系统由“仿真验证”转变为“实际应用”，具体内容如下所示。

\begin{enumerate}
    \item \textbf{完善硬件接入，形成完整的交互闭环}：之后会接入WEART TouchDIVER触觉手套和Oculus Rift S VR头显，使用WEART SDK和pyopenxr接口进行真实力反馈输出和连续手部mocap追踪，取代目前的键盘交互方式，解决非物理穿透和求解器发散的问题，完成人在回路验证，提高系统的现实感和临床适应性，回到原项目腹腔手术触觉训练的核心定位。
    
    \item \textbf{改进算法和仿真精度，提高真实感}：改进触觉反馈算法，根据真实生物组织的摩擦特性建立数据驱动的动态纹理生成模型，代替目前的预设纹理映射方式，提高纹理的细腻程度和材质的区分度；改进FEM求解器参数，使用隐式欧拉求解和Newton迭代方法，提高极端大变形场景下数值稳定性；改进腹腔器官粘弹性力学响应仿真，提高建模精度。
    
    \item \textbf{扩展系统的功能，提高系统的实用性}：扩展双手协同交互、多指精细控制和复杂手势识别的功能，适应更加复杂的手术操作环境；增加可交互的物体种类，完善场景配置，改进多物体接触检测算法，提高复杂场景下的系统稳定性；优化线程调度策略，进一步减小系统端到端延迟，提高视触同步精度。
    
    \item \textbf{推动场景落地，扩大应用范围}：把腹腔手术触觉训练当作主要的应用场景，加强系统与临床的契合度，给外科医生创建起标准的虚拟训练平台；而且拓宽应用范畴，让系统可以应用于康复训练、虚拟教学以及工业柔性装配等场景之中，并且凭借系统的算法通用性以及可扩展性来给出定制化的触觉仿真的方案。
\end{enumerate}

伴随着虚拟现实、人机交互以及智能硬件技术的发展，触觉反馈也成了下一代沉浸式交互的主要支持。本研究所完成的系统改造以及算法优化，给可变形物体触觉仿真提供了一个可以借鉴的硬件解耦方案和算法框架，之后再结合硬件接入以及功能优化，可以将触觉交互技术由理论仿真变为实际应用，给更加自然、真实、智能的人机交互提供技术支持。